\BourbakiExerciseGroup

\begin{enumerate}
\def\labelenumi{\arabic{enumi})}
\tightlist
\item
  A topology \(\mathfrak{G}\) on an ordered commutative group \(\mathbf{G}\) is said to be compatible with the ordered group structure of \(\mathbf{G}\) if it is compatible with the group structure of \(\mathbf{G}\) and if in addition the set \(\mathbf{G}_{+}\) of elements \(x \geqslant 0\) in \(\mathbf{G}\) is closed with respect to \(\mathfrak{G}\) .
\end{enumerate}

\begin{enumerate}
\def\labelenumi{\alph{enumi})}
\item
  Let \(\mathfrak{G}\) be a Hausdorff topology compatible with the ordered group structure of \(\mathbf{G}\) , and let \(\hat{\mathbf{G}}\) be the completion of \(\mathbf{G}\) (with respect to \(\mathfrak{G}\) ). In the group \(\hat{\mathbf{G}}\) , the closure \(\mathbf{P}\) of \(\mathbf{G}_{+}\) is such that \(y - x \in \mathbf{P}\) is a pre-ordering compatible with the group structure of \(\hat{\mathbf{G}}\) .
\item
  Let \(\theta\) be an irrational number, and consider the ordering on \(\mathbf{Q}^2\) for which the set of positive elements consists of (0, 0) and all pairs \((x, y)\) such that \(y - \theta x \geqslant 0\) . With respect to this ordering, \(\mathbf{Q}^2\) is a linearly ordered group, and the topology on \(\mathbf{Q}^2\) which is the product of the topology of the rational line by itself is compatible with this ordered group structure. But on \(\hat{\mathbf{G}} = \mathbb{R}^2\) , the relation \(y - x \in \mathbf{P} = \overline{\mathbf{G}}_+\) is not an ordering.
\item
  On a linearly ordered group G, the topology \(\mathfrak{C}_{0}(\mathbf{G})\) (Chapter I, § 2, Exercise 5) is the coarsest of all topologies compatible with the ordered group structure (distinguish two cases, according as the set of all \(x > 0\) has or has not a smallest element). Every topology on G which is compatible with the group structure of G and is finer than \(\mathfrak{C}_{0}(\mathbf{G})\) is compatible with the ordered group structure of G, and the mapping \(x \to x^{+}\) is continuous with respect to such a topology. This mapping is uniformly continuous with respect to \(\mathfrak{C}_{0}(\mathbf{G})\) , but not necessarily with respect to a group topology \(\mathfrak{C}\) strictly finer than \(\mathfrak{C}_{0}(\mathbf{G})\) {[}see b){]}.
\end{enumerate}

\begin{enumerate}
\def\labelenumi{\arabic{enumi})}
\setcounter{enumi}{1}
\tightlist
\item
  \begin{enumerate}
  \def\labelenumii{\alph{enumii})}
  \tightlist
  \item
    Let G be a lattice-ordered commutative group. For each \(x \geqslant 0\) in G, let \(\mathbf{I}(x)\) denote the interval \([-x, x]\) in G. A non-empty family \((c_{\alpha})\) of elements \(\geqslant 0\) in G is such that the intervals \(\mathbf{I}(c_{\alpha})\)
  \end{enumerate}
\end{enumerate}

form a fundamental system of neighbourhoods of 0 for a topology \(\mathfrak{G}\) on G compatible with the group structure of G if and only if the set of \(c_{\alpha}\) is directed (with respect to the relation \(\geqslant\) ) and for each \(\alpha\) there exists \(\beta\) such that \(2c_{\beta} \leqslant c_{\alpha}\) . If these conditions are satisfied, the mapping \(x \to x^{+}\) of G into G is uniformly continuous. \(\mathfrak{G}\) is Hausdorff if and only if \(\inf c_{\alpha} = 0\) ; and in this case \(\mathfrak{G}\) is compatible with the ordered group structure of G.

\begin{enumerate}
\def\labelenumi{\alph{enumi})}
\setcounter{enumi}{1}
\item
  On the group \(G = Z^{N}\) , the product of a countable infinity of linearly ordered groups Z, there is no non-discrete topology defined by the procedure of a), but the product of the discrete topologies on the factors Z is a topology on G, compatible with the ordered group structure of G, which is Hausdorff and not discrete.
\item
  On a linearly ordered group G the only Hausdorff topology defined by the procedure of a) is the topology \(\mathcal{G}_{0}(G)\) (Chapter I, § 2, Exercise 5).
\end{enumerate}

\begin{enumerate}
\def\labelenumi{\arabic{enumi})}
\setcounter{enumi}{2}
\tightlist
\item
  Let G be a lattice-ordered commutative group, and let M be the semigroup of major subsets of G. Every subset of M which is bounded above (resp. below) has a least upper bound (resp. greatest lower bound), and G can be canonically identified with a subset of M. Let G' denote the largest subgroup of M (the set of all symmetrizable elements of M).
\end{enumerate}

\begin{enumerate}
\def\labelenumi{\alph{enumi})}
\tightlist
\item
  Consider a topology \(\mathcal{G}\) on \(\mathbf{G}\) defined by the process of Exercise 2 \(a\) ); let \(\mathbf{V}_{\alpha}\) be the set of all pairs \((z, z')\) in \(\mathbf{M} \times \mathbf{M}\) such that
\end{enumerate}

\[
z - c _ {\alpha} \leqslant z ^ {\prime} \leqslant z + c _ {\alpha};
\]

show that the \(V_{\alpha}\) form a fundamental system of entourages of a uniformity U on M, and that the topology \(C'\) induced by U induces C on G. If C is Hausdorff, then U is a Hausdorff uniformity (remark that the elements of M, considered as subsets of G, are then closed), and the uniform space M thus defined is complete. Show that the mapping \((z, z') \to z + z'\) of M × M into M is uniformly continuous. Supposing that C is Hausdorff, show that in M the set of upper (resp. lower) bounds of any subset of M is closed with respect to \(C'\) (prove this first for the set of upper or lower bounds of an element of M, by considering the elements of M as subsets of D), and that the mapping

\[
(z, z ^ {\prime}) \rightarrow \inf (z, z ^ {\prime})
\]

of \(\mathbf{M} \times \mathbf{M}\) into \(\mathbf{M}\) is uniformly continuous (same method).

\begin{enumerate}
\def\labelenumi{\alph{enumi})}
\setcounter{enumi}{1}
\item
  Suppose from now on in this exercise that \(\mathfrak{G}\) is Hausdorff. Then the closure \(\overline{\mathbf{G}}\) of \(\mathbf{G}\) in \(\mathbf{M}\) is a lattice-ordered group which is complete in the topology induced by \(\mathfrak{G}'\) . The group \(\mathbf{G}'\) is closed in \(\mathbf{M}\) (remark that its closure is a subgroup of M) and the topology induced by \(\mathcal{G}'\) on \(\mathbf{G}'\) is compatible with the ordered group structure of \(\mathbf{G}'\) . For \(\mathbf{G}'\) to be equal to \(\overline{\mathbf{G}}\) it is sufficient that every non-empty open interval \(]a, b[\) in \(\mathbf{G}'\) should contain an element of \(\mathbf{G}\) , and this is always the case if \(\mathbf{G}\) is linearly ordered. In this case the topology induced on \(\mathbf{G}'\) by \(\mathcal{G}'\) is \(\mathcal{C}_0(\mathbf{G}')\) .
\item
  Take G to be the ordered group \(Q^{2}\) , the product of the linearly ordered group Q by itself; then \(G' = M = R^{2}\) with the product ordering. Show that there are only three distinct non-discrete Hausdorff topologies on G defined by the procedure of Exercise 2 a), and that they are obtained by taking the set of elements \(c_{\alpha}\) to be either the set of pairs \((x, y)\) such that x \textgreater{} 0 and y \textgreater{} 0, or the set of pairs \((x, 0)\) such that x \textgreater{} 0, or the set of pairs \((0, y)\) such that y \textgreater{} 0. For the first of these three topologies we have \(\overline{G} = G'\) , although there are non-empty open intervals of \(G'\) which contain no element of G; for the other two topologies, we have \(\overline{G} \neq G'\) . For each of these three topologies, the set of elements z \textgreater{} 0 in G is not open.
\item
  Take G to be the group \(Q^{2}\) endowed with the lexicographic order (Set Theory, Chapter III, § 2, no. 6); G is a non-Archimedean linearly ordered group. Then \(\overline{G}=G'\neq M\) ; M is linearly ordered but the topology \(C'\) on M is distinct from the topology \(\mathcal{C}_{0}(M)\) ; \(G'\) is isomorphic to the group \(R\times Q\) with the lexicographic order; in \(G'\) the subgroup \(H=R\times\{0\}\) is open and isolated, but on the quotient group \(G'/H\) the quotient topology is distinct from \(\mathcal{C}_{0}(G'/H)\) .
\end{enumerate}

4) a) Let E be a linearly ordered set. Let F be the Z-module of formal finite linear combinations of elements of E with coefficients in Z, and define a linearly ordered group structure on F by taking the set \(F_{+}\) of elements \(\geqslant 0\) to be the set consisting of 0 and all linear combinations \(\sum_{\xi \in E} n(\xi) \xi \neq 0\) such that \(n(\xi) > 0\) for the largest element \(\xi \in E\) such that \(n(\xi) \neq 0\) . Show that E can be identified with a cofinal subset of \(F_{+}\) .

\begin{enumerate}
\def\labelenumi{\alph{enumi})}
\setcounter{enumi}{1}
\item
  Suppose that E is well-ordered, and consider the group G of bounded mappings of E into F. Define a linearly ordered group structure on G by taking the set \(G_{+}\) of elements \(\geqslant0\) to be the set consisting of 0 and all bounded mappings x of E into F such that \(x(\xi)>0\) for the smallest \(\xi\in E\) such that \(x(\xi)\neq0\) (lexicographic ordering). Show that G is complete with respect to the topology \(\mathcal{C}_{0}(G)\) ; if moreover E is uncountable and is such that every segment \(\leftarrow,\xi]\) ( \(\xi\in E\) ) is countable, then every countable intersection of open sets in G is open, and every compact subset of G is finite (cf.~Chapter I, § 9, Exercise 4).
\item
  Suppose henceforth that E is well-ordered and uncountable but that every segment \(\leftarrow,\xi]\) is countable. For each \(\xi\in E\) , let \(c_{\xi}\) be the constant mapping of E into E whose value at every point is \(\xi\) . The \(c_{\xi}\) form a cofinal subset in \(G_{+}\) . Let \(E_{0}\) be a set obtained by adjoining an element \(\omega\) to E, and consider the topology G, on the set \(X=G\times E_{0}\) , generated by the following sets: (i) \(V_{a,b,\xi} = ]a\) , \(b[\times\{\xi\}\) for a \textless{} b in G and \(\xi\in E\) ; (ii) \(W_{a,b,\xi} -\) for a \textless{} b in G and \(\xi\in E\) such that \(c_{\xi} > \sup(|a|, |b|) -\) consisting of the pairs \((x, \omega)\) such that a \textless{} x \textless{} b, and the pairs \((x, \zeta)\) such that \(\zeta\in E\) , \(\zeta \geqslant \xi\) and either a \textless{} x \textless{} b or \(4c_{\zeta} + a < x < 4c_{\zeta} + b\) . Show that G is Hausdorff and that every compact subset of X (with respect to G) is finite. Moreover G operates continuously on X according to the law \((x, (y, \zeta)) \to (x + y, \zeta)\) ; but G does not operate properly on X, although conditions a), b), c) and d) of Chapter III, § 4, no. 2, Proposition 4 are satisfied and although, for each pair of compact subsets K, L of X, the set P(K, L) (Chapter III, § 4, no. 5, Theorem 1) is compact.
\end{enumerate}
% semantic-fragments: legacy-input-seam
\relax{}
